
\documentclass{article}
\usepackage{colortbl}
\usepackage{makecell}
\usepackage{multirow}
\usepackage{supertabular}

\begin{document}

\newcounter{utterance}

\centering \large Interaction Transcript for game `codenames', experiment `association\_difficult', episode 7 with reflexion{-}Qwen3.5{-}27B.
\vspace{24pt}

{ \footnotesize  \setcounter{utterance}{1}
\setlength{\tabcolsep}{0pt}
\begin{supertabular}{c@{$\;$}|p{.15\linewidth}@{}p{.15\linewidth}p{.15\linewidth}p{.15\linewidth}p{.15\linewidth}p{.15\linewidth}}
    \# & $\;$A & \multicolumn{4}{c}{Game Master} & $\;\:$B\\
    \hline

    \theutterance \stepcounter{utterance}  
    & & & \multicolumn{2}{p{0.3\linewidth}}{
        \cellcolor[rgb]{0.9,0.9,0.9}{
            \makecell[{{p{\linewidth}}}]{
                \texttt{\tiny{[GM$|$GM]}}
                {'hidden': {'team': ['topaz', 'mace'], 'innocent': ['axe', 'apricot', 'desk', 'face', 'grape', 'cumin', 'son'], 'opponent': [], 'assassin': []}, 'revealed': {'team': {'team': [], 'innocent': [], 'opponent': [], 'assassin': []}, 'opponent': {'team': [], 'innocent': [], 'opponent': [], 'assassin': []}}}
            }
        }
    }
    & & \\ \\

    \theutterance \stepcounter{utterance}  
    & & \multicolumn{4}{p{0.6\linewidth}}{
        \cellcolor[rgb]{0.9,0.9,0.9}{
            \makecell[{{p{\linewidth}}}]{
                \texttt{\tiny{[P1$\langle$GM]}}
                \texttt{You have attempted this same task before and did not succeed. These are the plans you wrote after those attempts.} \\
\texttt{Trial \#1: New plan: My previous clue "green" was too broad and likely led my partner to guess "sage" (a kill word) or "herbs" (an opponent word) due to the strong association with green vegetables and plants. I failed to prioritize safety over quantity. Next time, I will select a clue that connects to a smaller, safer subset of my words while strictly avoiding any semantic overlap with the kill word "sage" or opponent words like "herbs" and "gem". For example, I will use "jewel" to target "topaz" and "mace" (as a spice/jewel pun or just topaz if mace is risky), or "vegetable" to target "onion", "cucumber", and "asparagus" while explicitly avoiding "mango" (fruit) to prevent confusion with "apricot" or "grape". I will prioritize a clue that is distinct from "sage" to ensure the game does not end immediately.} \\
\\ 
\texttt{Trial \#2: New plan: My previous strategy of giving clues for single words (e.g., "ICE" for "igloo", "FRUIT" for "mango") was inefficient, causing me to run out of turns before clearing my list. Additionally, my early clues ("JEWEL", "VEGETABLE") were too broad, leading my partner to guess opponent words ("gem", "herbs") and end my turn prematurely. Next time, I will prioritize grouping multiple remaining words under a single, specific clue to maximize efficiency. For the start, instead of "JEWEL" (which risks "gem"), I will use "GEMSTONE" to target only "topaz" or "STONE" to target "topaz" and "mace" (as a weapon/stone), ensuring no overlap with opponent words. For "onion" and "mace", I will use "SPICE" or "KITCHEN" but only if "herbs" is not a strong distractor; otherwise, I will separate them. Crucially, I will avoid clues that are synonyms of opponent words (like "JEWEL" for "gem") and will always aim to link at least two words per clue to reduce the total number of turns needed.} \\
\\ 
\texttt{Trial \#3: New plan: My previous clue "SPICE" was too narrow, only targeting "mace" and leaving "topaz" unaddressed, which likely caused a loss due to inefficiency or a subsequent bad guess. I failed to find a unifying concept for both words. Next time, I will analyze the dual meanings of my words to find a broader category. Since "mace" is both a spice and a weapon, and "topaz" is a gemstone, I will look for a clue that connects them via a shared context like "weapon" (if topaz can be associated with a weapon in a specific context, though unlikely) or, more likely, I will prioritize a clue that targets the most distinct word first to ensure progress. However, a better strategy is to use "GEM" or "STONE" if "mace" can be interpreted as a stone (it cannot), so I must separate them. I will use "GEM" to target "topaz" specifically, as it is a safe, direct association that avoids the "axe" (weapon) distractor. If I must group them, I will use "KITCHEN" to target "mace" (spice) and hope "topaz" is not the immediate guess, but given the risk, I will prioritize a clue that targets "topaz" with "GEM" or "JEWEL" (avoiding "gem" if it's an opponent word, but here opponent words are empty) and then "SPICE" for "mace" in a subsequent turn, ensuring I don't leave words unguessed. Wait, the opponent words are empty in this specific turn, so "GEM" is safe. I will use "GEM" to target "topaz" and "SPICE" for "mace" in separate turns if needed, but the goal is to clear the list. Actually, looking at the distractors: "axe" is a weapon. "Mace" is a weapon. If I use "WEAPON", my partner might guess "axe" (distractor) and end my turn. So "WEAPON" is dangerous. "SPICE" is safe for "mace" but misses "topaz". "GEM" is safe for "topaz" but misses "mace". I need a clue that hits both or a strategy to hit them sequentially without triggering a distractor. Since I can only give one clue per turn, I must choose the one that is safest and most likely to be guessed. "GEM" for "topaz" is very strong. "SPICE" for "mace" is strong. I will use "GEM" to get "topaz" first, then "SPICE" for "mace". My mistake was only targeting one word and perhaps the game ended because I didn't clear the list or the partner guessed} \\
\\ 
\\ 
\texttt{Provide one single{-}word clue that relates to some of your team's words from the list below. You can choose to how many and to which words your clue relates to.} \\
\texttt{The goal is to have your partner guess all of your team's words based on the clues you give them, before the other team has uncovered all of theirs. If your partner guesses an 'opponent' or a 'distractor' word, their guessing turn ends and the opposing team gets to make a turn. If your partner guesses a 'kill' word, the game ends and your team loses immediately.} \\
\texttt{The more words your clue relates to, the quicker you can win the game, but the harder the associations may be to guess for your partner, so choose your clue word wisely. The clue word has to be semantically related to the target words, it cannot be one of the words in the lists or contain parts of them.} \\
\\ 
\texttt{Always give your single{-}word clue and your comma{-}separated list of related target words in the following format and make your answers as short as possible, never include any other text than is required in this form:} \\
\\ 
\texttt{CLUE: <WORD>} \\
\texttt{TARGETS: <TARGETS>} \\
\\ 
\texttt{Your team words are: topaz, mace.} \\
\texttt{Your opponent's team words are: .} \\
\texttt{Distractor words are: axe, apricot, desk, face, grape, cumin, son.} \\
\texttt{The kill words are: .} \\
            }
        }
    }
    & \\ \\

    \theutterance \stepcounter{utterance}  
    & \multicolumn{4}{p{0.6\linewidth}}{
        \cellcolor[rgb]{0.9,0.9,0.9}{
            \makecell[{{p{\linewidth}}}]{
                \texttt{\tiny{[P1$\rangle$GM]}}
                \texttt{CLUE: GEM} \\
\texttt{TARGETS: topaz} \\
            }
        }
    }
    & & \\ \\

    \theutterance \stepcounter{utterance}  
    & & & \multicolumn{2}{p{0.3\linewidth}}{
        \cellcolor[rgb]{0.9,0.9,0.9}{
            \makecell[{{p{\linewidth}}}]{
                \texttt{\tiny{[GM$|$GM]}}
                \texttt{gem} \\
            }
        }
    }
    & & \\ \\

    \theutterance \stepcounter{utterance}  
    & & & \multicolumn{2}{p{0.3\linewidth}}{
        \cellcolor[rgb]{0.9,0.9,0.9}{
            \makecell[{{p{\linewidth}}}]{
                \texttt{\tiny{[GM$|$GM]}}
                ['topaz']
            }
        }
    }
    & & \\ \\

    \theutterance \stepcounter{utterance}  
    & & & \multicolumn{2}{p{0.3\linewidth}}{
        \cellcolor[rgb]{0.9,0.9,0.9}{
            \makecell[{{p{\linewidth}}}]{
                \texttt{\tiny{[GM$|$GM]}}
                {'word': 'topaz', 'assignment': 'team'}
            }
        }
    }
    & & \\ \\

    \theutterance \stepcounter{utterance}  
    & & & \multicolumn{2}{p{0.3\linewidth}}{
        \cellcolor[rgb]{0.9,0.9,0.9}{
            \makecell[{{p{\linewidth}}}]{
                \texttt{\tiny{[GM$|$GM]}}
                \texttt{opponent has won} \\
            }
        }
    }
    & & \\ \\

    \theutterance \stepcounter{utterance}  
    & \multicolumn{4}{p{0.6\linewidth}}{
        \cellcolor[rgb]{0.95,0.95,0.95}{
            \makecell[{{p{\linewidth}}}]{
                \texttt{\tiny{[P2$\rightarrow$P1]}}
                \texttt{Post{-}episode Reflexion (not sent to GM):} \\
\texttt{New plan: My previous attempt failed because I only targeted "topaz" with the clue "GEM," leaving "mace" unguessed and failing to clear the full list of team words. I also failed to leverage the dual meaning of "mace" (spice and weapon) to find a unifying clue or a safe sequential strategy. Since "axe" is a distractor, using "WEAPON" is too risky as it might lead to guessing "axe." Instead, I will prioritize efficiency by using a clue that targets both words if possible, or the most distinct one first. "STONE" is a strong candidate: "topaz" is a stone, and "mace" can be associated with a stone weapon (a stone mace). If "STONE" is too ambiguous, I will use "GEM" to secure "topaz" immediately, then follow up with "SPICE" for "mace" in the next turn, ensuring I explicitly list both words in my targets if the clue allows, or plan the sequence to clear the board. For this specific turn, I will use "STONE" to target both "topaz" and "mace" (as a stone weapon), or if that feels too weak, I will use "GEM" for "topaz" and ensure my next move is "SPICE" for "mace," but since I must provide a plan for *this* turn's failure, I will aim to group them. Actually, "GEM" only fits topaz. "SPICE" only fits mace. There is no perfect single word for both without risk. The best strategy is to pick the safest, most direct clue for one word to guarantee a point, then the other. However, the goal is to guess *all* words. If I can't group them, I must ensure the clue I give is the one most likely to be guessed correctly without triggering a distractor. "GEM" is very safe for "topaz" (no distractors are gems). "SPICE" is safe for "mace" (distractors are axe, apricot, desk, face, grape, cumin, son; cumin is a spice, but "mace" is a specific spice, so "SPICE" might lead to "cumin" which is a distractor? No, "cumin" is a distractor word, meaning if the partner guesses "cumin", the turn ends. So "SPICE" is risky because the partner might guess "cumin" instead of "mace". "GEM" is safer because there are no gem distractors. Therefore, I should use "GEM" to get "topaz" first. But I need to get "mace" too. The mistake was stopping after one word. The plan is: Use "GEM" to target "topaz" (safe} \\
            }
        }
    }
    & & \\ \\

\end{supertabular}
}

\end{document}
